\section{Evaluation}\label{sec:offline-evaluation}

We establish the practicality of the scheme by integrating the protocol into a
mobile wallet and benchmarking it on both iOS and Android. Of the six
participants of~\Cref{sec:offline-definitions}, only the user's wallet is
resource constrained, so we evaluate the mobile phone and the secure element
that make it up. The intermediary, the central bank, the registration authority,
and the auditor are expected to run on servers. The phone and the secure element
take four roles across the protocol: the holder who receives funds during a
withdrawal, the payer and the payee who exchange possession of a token during a
payment, and the depositor who redeems it.

\subsection{Implementation}

\paragraph{Secure element.} The choice of secure element determines what a
benchmark measures. We discard three options. Software emulation would ignore
both the constrained environment of a real device and the cost of communicating
with it. TEE based secure elements are unsuitable for the reasons given
in~\Cref{sec:offline-construction}. The HSM based secure elements integrated
into the Apple and Android platforms are not extensible and do not implement our
scheme. We therefore use an external secure element, realised as a Javacard
applet on a smart card, which protects the key and the logic together and
supports the elliptic curve operations the protocol requires~\cite{javacard}.

The applet exposes two functions, receive and sign. Recall that
the payer's card contributes a signature of knowledge over the payer's and the
payee's commitments, $\com_i$ and $\com_{i+1}$. The card signs only triples $(v,
\com_i, \com_{i+1})$ in which $\com_i$ is the commitment it generated when it
received the token, and, once that token is spent, only against the single
destination $\com_{i+1}$ it first signed. The first constraint prevents a token
that has already been passed on from being spent again, and the second prevents
the current token from being sent to two recipients. It is this state machine,
rather than any property of the cryptography, that realises the single-use
enforcement the construction attributes to the secure element. Repeated
invocations of sign on an unchanged $\com_{i+1}$ are permitted, so that
errors on the NFC channel, which are not uncommon, do not cost the user a
payment. Invoking receive deletes the stored commitment and generates a
new one, which frees the slot but destroys any token held in it, so a device
that resets a slot to escape the single destination rule loses the token in
doing so. A single card holds several tokens by partitioning this state into
slots, with each message carrying a slot identifier.

\paragraph{Mobile wallet.} The wallet is an application that withdraws and
transfers tokens, communicating with the smart card over NFC, with the payee
over BLE, and with the intermediary and the registration authority, both exposed
as cloud services, over HTTPS.

\subsection{Protocols}

Withdrawal consists of local computation on the phone, an NFC exchange with the
card to obtain $s_0$ and $\com_0$, and a single HTTPS call to the intermediary,
which contacts the central bank on the user's behalf.
Figure~\ref{fig:offline/withdraw-sequence} gives the message sequence.

A payment is preceded by an out of band negotiation of its terms. The payee
displays a QR code carrying the payment terms, a BLE service identifier, and the
material for a Diffie-Hellman exchange, and advertises a BLE GATT service with a
Notify characteristic for data sent to the payer and a Write without response
characteristic for data received from them. The payer scans the code, completes
the exchange and confirms the terms, after which every message is encrypted and
CBOR encoded. Figure~\ref{fig:offline/payment-sequence} gives the message
sequence for the payment itself.

\paragraph{Interrupted payments.} A payment that stops part way leaves the two
wallets in one of two states, and neither loses value. If the payer's card has
signed but the payee never received the payment, the card remains locked to the
destination it signed, and the payer can retry against that same destination
until it succeeds; the token is unspendable elsewhere in the meantime, and is
recovered by completing the payment rather than by any repair step. If the payee
stored the token but the acknowledgement back to the payer was lost, the payer's
card still holds $\com_i$ and has not deleted it, but it will sign only the
destination it already signed, so it cannot spend the token a second time.  The
payee holds a token that verifies and can spend or deposit it. In both cases the
resolution is a retry over NFC or BLE.

\subsection{Results}

We benchmark on a Samsung Galaxy A34 running Android 14 and an iPhone 13 mini
running iOS 17.6.1, both using a cryptographic library built with Go 1.21.8.
The scheme is instantiated over the BLS12-381 pairing system, at roughly $128$
bits of security; the construction relies only on the bilinearity of the
pairing, so other pairing systems would serve. The secure element is an NXP 1ID
white plastic card with a Smart MX D600 chip offering $400\,\mathrm{KiB}$ of
available memory, running JCOP 4.5 with NXP's JCOPx extensions at version 1.1.4,
and communicating with the phone over NFC. We benchmark the protocol both with
and without the verifiable encryption of the payer's identifier, to separate the
cost of auditing from the cost of the payment itself.

Table~\ref{tab:offline-timings} reports the cost of each operation on both
platforms. Withdrawal and the payer side of a payment each complete in under
half a second, of which the smart card accounts for under $400\,\mathrm{ms}$.
Auditing adds no measurable overhead for the payer, whose audit ciphertext and
proof are computed once regardless of how long the chain already is, and are
folded into the transfer signature as described
in~\Cref{sec:offline-instantiation}.

Table~\ref{tab:offline-scaling} reports the payee's cost as the chain grows.
Verification is linear in the number of hops, since each hop contributes one
signature of knowledge and one audit proof to check, and a chain of $32$ hops
still verifies in under a second. Auditing costs the payee a further $10$ to
$20\%$, one additional proof per hop. The transcript $\algo{hist}_n$ grows
linearly in the same way, as it does in the transferable ecash schemes
of~\Cref{sec:offline-related} that identify only the double spender.

Table~\ref{tab:offline-storage} reports the resources the protocol asks of the
smart card. The applet and its per slot state fit comfortably within the card's
memory, leaving room for a user to hold several tokens at once.

\begin{table}[t] \centering \small \begin{tabular}{l|rr|rr} \toprule &
\multicolumn{2}{c|}{Android} & \multicolumn{2}{c}{iOS} \\ 
%
Operation & no audit & audit & no audit & audit \\
%
 \midrule Withdrawal, device            & & & & \\ 
% FILL
 Withdrawal, smart card        & & & & \\ 
% FILL
 Payment, payer device & & & & \\ 
% FILL 
Payment, payer smart card     & & & & \\
 % FILL 
 Payment, payee device, 1 hop  & & & & \\ 
% FILL 
Deposit, device               & & & & \\ 
% FILL
Deposit, smart card           & & & & \\ 
% FILL 
\bottomrule \end{tabular} \caption{Wall clock cost in milliseconds of each
protocol operation, with and without the verifiable encryption of the payer's
identifier for the auditor.  Smart card figures include the NFC round trip. Each
figure is the mean of $[\,\cdot\,]$ runs, with the standard deviation in
parentheses.} \label{tab:offline-timings} \end{table}

\begin{table}[t] \centering \small \begin{tabular}{l|rrrrrr} \toprule Hops $n$ &
1 & 2 & 4 & 8 & 16 & 32 \\ \midrule Payee verification, no audit (ms) & & & & &
& \\ 
% FILL
 Payee verification, audit (ms)    & & & & & & \\
  % FILL 
  Transcript $|\algo{hist}_n|$ (B)  & & & & & & \\ 
% FILL
 \bottomrule \end{tabular} \caption{Scaling of the payee's verification cost and
of the token transcript with the number of hops the token has taken since
withdrawal. Verification figures are for Android, and iOS follows the same
trend. Each figure is the mean of $[\,\cdot\,]$ runs.}
\label{tab:offline-scaling} \end{table}

\begin{table}[t] \centering \small \begin{tabular}{l|r} \toprule Quantity &
Value \\ \midrule Applet code size  & \\ 
% FILL 
Persistent state per slot & \\
 % FILL 
 Slots available in $400\,\mathrm{KiB}$ & \\ 
% FILL 
Payment message, payer to payee & \\ 
% FILL 
Per hop increment to $\algo{hist}$ & \\
 % FILL 
 \bottomrule \end{tabular} \caption{Storage and communication requirements on
the smart card and the NFC and BLE channels. A slot holds the state for one
token.} \label{tab:offline-storage} \end{table}

These figures sit comfortably within the latency expected of an everyday
payment, on hardware that is already deployed. The comparison that matters here
is not with the ecash schemes of~\Cref{sec:offline-related}, which are not
benchmarked on constrained hardware, but with the alternative of putting the
whole protocol on the secure element. A smart card cannot carry the pairing
operations that the proofs of~\Cref{sec:offline-instantiation} require, so
without the split the scheme would not run at all. The split is what makes the
cost on the card comparable to that of an ordinary digital signature.

\paragraph{State storage requirements.} The central bank stores two lists, of
issued and of deposited tokens. The list of deposited tokens is needed to detect
double spending and cannot be pruned outright, though the epoch based expiry
of~\Cref{sec:offline-construction} bounds a token's validity and so allows
entries to be pruned once that span elapses. The list of issued tokens supports
monetary policy and is not required by the protocol. Intermediaries store user
accounts to determine reserve requirements, again not a protocol requirement,
and users store the transcripts of any unspent tokens they hold.
